\chapter{问题设定与假设}\label{sec:setup}

投资者的连续决策构成一系列有限时域优化问题。每期开始时，已观测的市场信息、财富、持仓和参考点构成当前情境；在这一情境下，给定策略产生一组可能的未来投资路径。本文以这些路径终端相对结果的分布定义条件 CPT 目标，并用投影残差衡量策略对当期一阶最优性条件的满足程度。市场与投资状态的变化由此进入同一在线优化问题。

\section{非平稳有限时域投资组合马尔可夫决策过程}

将一次有限时域决策称为一个时期（episode），编号为 $k=1,\ldots,K$，每个时期包含固定的 $h\geq 1$ 个决策步。第一个时期从 $t_1=1$ 开始，此后满足 $t_{k+1}=t_k+h$。本文用 $u$ 表示日历时点；实验中的上涨、震荡等市场阶段可以跨越多个时期，也可以在时期内部切换。一个时期内部的状态和投资组合动作分别写为
\[
 s_u=(m_u,X_u,r_u,\omega_{u-1}),\qquad
 \omega_u\in\Delta_N:=\{\omega\in\R_+^{N+1}:\textstyle\sum_{i=0}^N\omega_i=1\}.
\]
其中，$m_u$ 包含市场因子，$X_u$ 表示账户财富与固定初始财富之比，$r_u$ 是采用相同归一化单位表示的参考点，资产 $0$ 表示现金。令现金收益率 $R_{0,u+1}=0$，交易成本系数 $c_{\rm tc}\geq 0$，则财富递推为
\begin{equation}\label{eq:wealth-update}
 X_{u+1}=X_u\left(1+\omega_u^\top R_{u+1}
 -c_{\rm tc}\norm{\omega_u-\omega_{u-1}}_1\right),\qquad X_1=1.
\end{equation}
式中，$\omega_u^\top R_{u+1}$ 是该期组合收益，$\norm{\omega_u-\omega_{u-1}}_1$ 衡量相邻两期组合权重的变动。市场模型包含条件收益与状态转移核，该转移核可以随时期 $k$ 以及时期内时点变化。在给定状态和动作后，转移核不直接依赖策略参数。所有随机变量和条件转移核均可测，状态空间与动作空间均为标准 Borel 空间。初始参考点非负。允许的收益与交易成本应使式~\eqref{eq:wealth-update} 中的财富总增长因子严格为正，并满足后文给出的有界分析域。

在时期 $k$ 开始时，信息集 $\cF_k$ 包含已经观测到的历史、当前参数、初始状态以及已经拟合的模拟器。记
\[
 z_k=(m_{t_k},X_{t_k},\omega_{t_k-1}).
\]
本文仅将 $r_k:=r_{t_k}$ 用作时期起点参考点的简写；按日历时间记录的参考点仍使用下标 $u$ 或明确日期 $t_k+j$。对 $0\leq j\leq h$，定义时期内记号 $r_{k,j}:=r_{t_k+j}$ 与 $X_{k,j}:=X_{t_k+j}$。因此，$r_k=r_{k,0}$，且 $r_{k+1}=r_{k,h}$。真实目标与模拟目标分别为
\begin{equation}\label{eq:conditional-targets}
 J_k(\theta,r_{k})=J(P_k,z_k,r_k;\theta),\qquad
 \widetilde J_k(\theta,r_{k})=J(\widetilde P_k,z_k,r_k;\theta).
\end{equation}
全文对 $J_k(\theta,r_k)$ 和 $\widetilde J_k(\theta,r_k)$ 求梯度时，只对 $\theta$ 求导，并固定 $r_k$ 与 $z_k$。学习样本生成前，模拟器和初始状态也保持冻结。下文所有条件期望均表示在这些信息给定时的正则条件期望。$\theta_k$、$M_k$ 以及选定的截断层级均为 $\cF_k$ 可测。新时期可以纳入新观测数据重新拟合模拟器，但时期 $k$ 的模拟器一旦冻结，在该时期采样期间便不再改变。这样的条件化处理明确区分了时期之间的信息更新与时期内部的梯度采样。

冻结模拟器意味着在本期学习采样中保持拟合规则与参数不变。该规则仍可包含不同日期的转移核、预先给定的市场阶段以及状态依赖的收益分布。因此，有限时域内部允许非平稳变化；模拟器与真实过程之间的差异由后文的模拟梯度误差刻画。

目标的跨期变化有不同来源。外生非平稳指市场收益分布或条件转移规律本身随日历时间变化，例如市场冲击使同一状态和动作下的收益规律发生改变。另一方面，过去的配置决定已实现财富和持仓，财富经历又决定参考点；这些内生形成的起始信息通过 $z_k,r_k$ 改变当期条件目标。这些状态由此前的策略与市场实现共同生成，使目标变化具有内生来源。即使市场转移规律保持不变，从不同财富、持仓或参考点出发，也可能面对不同的条件目标。将财富和参考点纳入状态后，固定的递推规则仍可构成时间齐次的状态转移。因此，状态随路径演化与转移核本身随时间变化需要分别判断，二者对目标的共同影响由后文的变化预算记录。

\section{非对称动态参考点}

参考点表示投资者评价财富时使用的比较水平。例如，账户价值从初始 100 万元增长到 120 万元时，归一化财富为 1.2；若经历上涨后形成的参考点为 1.15，则当前相对结果为 0.05。该量始终以初始财富为单位，保留了此前财富积累与参考适应的共同影响。现金与股票都属于总账户，现金收益设为零，现金比例仍会改变财富分布和未来风险。

固定 $\eta_+,\eta_-\in[0,1]$，定义参考点更新映射
\begin{equation}\label{eq:reference-update}
 \Phi(r,x)=
 \begin{cases}(1-\eta_+)r+\eta_+x,&x\ge r,\\
 (1-\eta_-)r+\eta_-x,&x<r.
 \end{cases}
 \qquad r_{u+1}=\Phi(r_u,X_{u+1}).
\end{equation}
一种常见的行为设定是 $\eta_+>\eta_-$，即参考点在财富高于原参考点时调整得更快。这一设定与盈利后适应较快的行为动机一致\citep{arkes2008reference}。本文将两侧适应率作为给定的模型参数，优化定理适用于所声明区间内的取值；个体参数的校准属于行为实验的测量内容。

每条模拟轨迹都必须沿自身的财富序列递归生成对应参考点，并据此定义时期末相对结果
\begin{equation}\label{eq:relative-outcome}
 Y_k(\tau)=X_{k,h}(\tau)-r_{k,h}(\tau),
\end{equation}
其中，$X_{k,h}(\tau)=X_{t_k+h}(\tau)$ 与 $r_{k,h}(\tau)=r_{t_k+h}(\tau)$ 分别是轨迹 $\tau$ 的期末财富和期末参考点。每条轨迹均从已保存的时期 $k$ 起点信息初始化。式~\eqref{eq:wealth-update} 中的 $X_1=1$ 只适用于整个实验的最初时点，后续时期继承上一时期结束时的财富。若将真实执行路径上观测到的同一个期末参考点用于所有模拟轨迹，模拟轨迹的相对结果会随之改变，进而定义出另一个目标函数。

下文将 $Y_k$ 统一称为终端相对结果，其单位与归一化财富相同。整个账户自实验开始的累计收益率为归一化财富减 1，而 $Y_k$ 减去的是随路径变化的参考点，两者分别用于投资表现与 CPT 评价。动态参考点包括对称适应和非对称适应：前者使用相同的两侧适应率，后者允许二者不同。静态参考点则在执行路径和所有模拟路径中均令 $\eta_+=\eta_-=0$。动态描述递推机制，即使某一步财富恰好等于参考点，该机制仍属于动态设定。

\section{正则化累积前景目标}

一条轨迹提供一个终端相对结果，多条可能轨迹共同定义一个随机变量。CPT 对这一随机变量的分布进行评价：先计算每个结果在收益侧或损失侧的效用幅度，再考察超过各效用水平的概率，最后分别加权累积两侧贡献。这样，收益大小与其出现机会共同影响目标。有限样本实现中的排序用于表示经验分布，属于总体目标的估计步骤。

记 $y^+=\max(y,0)$、$y^-=\max(-y,0)$。固定 $\eps_v>0$，对 $x\geq0$ 定义
\begin{equation}\label{eq:value-functions}
 u_{\alpha,\eps_v}(x)=(x^2+\eps_v^2)^{\alpha/2}-\eps_v^\alpha,
 \quad v_+(y)=u_{\alpha_+,\eps_v}(y^+),\quad
 v_-(y)=\lambda u_{\alpha_-,\eps_v}(y^-),
\end{equation}
其中 $0<\alpha_+,\alpha_-\leq1$，$\lambda>1$。收益侧与损失侧的效用幅度在零点均为零，并且作为 $y$ 的函数在包括零点在内的定义域上连续可微；本文不要求它们关于 $y$ 具有全局二阶导数。两侧的全局 Lipschitz 常数可以分别取为
\begin{equation}\label{eq:value-lipschitz}
 L_v^+=\alpha_+\eps_v^{\alpha_+-1},\qquad
 L_v^-=\lambda\alpha_-\eps_v^{\alpha_--1}.
\end{equation}
这里的平滑是价值函数设定的一部分，会影响所有正效用幅度。经典幂价值函数在收益侧的凹性与损失侧的凸性，也不能直接沿用到这一平滑形式：在零点附近，平滑后的收益侧变换具有局部凸性，对应的带符号损失侧具有局部凹性。本文使用的是式~\eqref{eq:value-functions} 明确定义的正则化模型，后续证明依赖其可微性、Lipschitz 性及分析域内的效用界。固定平滑参数时所得结论均针对这一目标；逼近未平滑 CPT 的极限还需要单独分析。

给定严格递增且归一化分母为正的基础权重函数 $w_\pm$，定义
\begin{equation}\label{eq:prob-weights}
 w_\pm^\eps(p)=
 \frac{w_\pm(\eps+(1-2\eps)p)-w_\pm(\eps)}
 {w_\pm(1-\eps)-w_\pm(\eps)},\qquad 0<\eps<\tfrac12.
\end{equation}
本文的行为设定采用 Tversky 与 Kahneman 提出的函数族
\[
 w_\pm(p)=\frac{p^{\beta_\pm}}
 {\{p^{\beta_\pm}+(1-p)^{\beta_\pm}\}^{1/\beta_\pm}},
\]
并选择参数，使权重函数在正则化区间上递增。理论结果同样适用于满足上述条件的其他基础权重函数；受控数值验证采用 $w(p)=p^3$。由归一化定义可得 $w_\pm^\eps(0)=0$ 和 $w_\pm^\eps(1)=1$。所有价值参数与权重参数在不同样本量和多层层级之间保持固定。当收益侧与损失侧采用不同权重函数时，归一化可能改变两侧的相对尺度，这种尺度变化属于模型设定本身。

在真实轨迹分布或模拟轨迹分布下，记收益侧与损失侧效用超过阈值的尾概率（即生存函数）为
\[
 S_{k,\theta}^{\pm}(z)=\Pp_{k,\theta}\{v_\pm(Y_k)>z\}.
\]
若 $0\leq v_\pm(Y_k)\leq B_v$，则对应的时期目标为
\begin{equation}\label{eq:cpt-objective}
 J_k(\theta,r_{k})=\int_0^{B_v}w_+^\eps(S_{k,\theta}^+(z))\,dz
 -\int_0^{B_v}w_-^\eps(S_{k,\theta}^-(z))\,dz.
\end{equation}
模拟分布下的生存概率相应定义 $\widetilde J_k$。全文统一采用严格不等式定义生存概率，因此允许效用分布存在原子，也允许样本出现并列值。特别地，$|J_k|,|\widetilde J_k|\leq B_J:=B_v$。对一条固定轨迹而言，其效用泛函不显式依赖 $\theta$；策略参数只通过轨迹的概率分布影响目标。

\section{移动驻点集合与动态局部遗憾}

在第 $k$ 期，数学优化问题是对 $\theta\in\Theta$ 最大化 $J_k(\theta,r_k)$。策略在这一期内保持同一组参数，下一期则面对新的条件目标。一般非凸目标的全局最大点难以直接求得，本文先考察可行参数扰动是否还具有一阶改进空间。若参数位于边界，即使原始梯度不为零，其指向也可能全部被约束排除，因此评价时需要同时考虑目标梯度和参数域。

记 $g_k=\nabla J_k(\theta_k,r_{k})$、$\widetilde g_k=\nabla\widetilde J_k(\theta_k,r_{k})$。固定 $a_0>0$ 和 $\gamma>0$，定义归一化方向映射与归一化投影残差映射，后者在下文简称投影残差：
\begin{equation}\label{eq:normalized-map}
 \mathsf d(g)=\frac{g}{\max\{\norm{g}_2,a_0\}},\qquad
 Q_k(\theta)=\frac{\Pi_\Theta(\theta+\gamma \mathsf d(\nabla J_k(\theta,r_{k})))-\theta}{\gamma}.
\end{equation}
其中，阈值 $a_0$ 与现金资产对应的 Dirichlet 浓度参数 $a_{0,u}$ 是两个不同对象。受约束的一阶驻点集合定义为
\begin{equation}\label{eq:stationary-set}
 \Sta_k=\{\theta\in\Theta:\ip{\nabla J_k(\theta,r_{k}),u-\theta}\le0
 \text{ for every }u\in\Theta\}=\{\theta\in\Theta:Q_k(\theta)=0\}.
\end{equation}
上述两个集合相等，可由投影变分不等式以及梯度的正比例缩放得到。连续目标函数在紧集 $\Theta$ 上能够取得最大值，其最大化点满足该一阶条件，因此 $\Sta_k$ 非空。$Q_k$ 的连续性进一步保证 $\Sta_k$ 为闭集。该集合也可能包含鞍点和非全局驻点，所以后续结论针对受约束一阶驻点，而不把它们统一解释为全局最优解。

梯度范数平方 $\norm{g_k}_2^2$ 衡量目标关于参数的一阶敏感程度，投影残差平方 $\norm{Q_k(\theta_k)}_2^2$ 还包含方向归一化与参数约束的作用。若候选更新点位于 $\Theta$ 内，投影不改变该点，此时 $Q_k(\theta_k)=g_k/\max\{\norm{g_k}_2,a_0\}$；再满足 $\norm{g_k}_2\le a_0$ 时，才有 $\norm{Q_k(\theta_k)}_2^2=\norm{g_k}_2^2/a_0^2$。因此，比较两类曲线时必须保留归一化阈值和触边信息。梯度估计 $\widehat g_k$ 的范数平方还含有采样误差，其时间曲线需要与总体梯度的性质分开解释。

本文使用的主要评价准则是逐时期真实投影残差平方的平均值
\begin{equation}\label{eq:primary-residual}
 \mathcal E_K=\frac1K\sum_{k=1}^K\E\norm{Q_k(\theta_k)}_2^2.
\end{equation}
进一步参考 Aydore 等在各期自身参数处加权历史梯度的思路\citep{aydore2019dynamic}，固定窗口长度 $w\geq1$、衰减系数 $0<\rho\leq1$，并令
$A_k=\sum_{j=0}^{\min(w-1,k-1)}\rho^j$，定义加权窗口诊断量
\begin{equation}\label{eq:smoothed-gradient}
 \overline Q_k=A_k^{-1}\sum_{j=0}^{\min(w-1,k-1)}\rho^j
 Q_{k-j}(\theta_{k-j}),
\end{equation}
据此定义投影残差形式的动态局部遗憾，下文未作特别说明时均指这一量：
\begin{equation}\label{eq:dynamic-regret}
 \Reg_{\rho,w}^{Q}(K)=\sum_{k=1}^K\norm{\overline Q_k}_2^2.
\end{equation}
相较于 Aydore 等的定义，式~\eqref{eq:dynamic-regret}作了两项明确调整。其一，将无约束梯度替换为与本文归一化投影更新相对应的 $Q_k$；其二，在最初不足 $w$ 期时，只对已经出现的时期按 $A_k$ 归一化。原文使用完整窗口的固定权重和，并把初始缺少的历史项设为零；两种归一化在窗口填满后相同。这些调整保留了在历史自身参数处评价的思路，其上界由第~\ref{sec:results}章单独建立。

计算式~\eqref{eq:dynamic-regret}时，先对近期残差向量加权平均，再取范数平方，最后对所有时期累加。直接把各期梯度范数平方相加，会得到另一个累计一阶指标；先对范数平方做滑动平均，也会失去向量之间的方向信息。只有窗口宽度为 1 时，动态局部遗憾才退化为逐期投影残差平方的累计和。不同时期的向量可能抵消，故较小的窗口量需要结合式~\eqref{eq:primary-residual}共同判断。

通常的动态遗憾累计的是各期目标值与该期比较策略目标值之间的差额；取逐期最优策略为比较对象时，对应最优值差额。本文采用的一阶局部准则与这种目标值差额具有不同含义，到驻点集合的参数距离则另由假设~\ref{ass:error-bound}联系起来。式~\eqref{eq:dynamic-regret}本身是随机累计量，理论上界对其取期望；实验中除以已经运行的时期数得到平均动态局部遗憾。所谓次线性增长，是指 $\E\Reg_{\rho,w}^Q(K)=o(K)$，也就是该期望累计量除以 $K$ 后趋于零。这一增长性质表示单位时期的累计一阶偏离趋于零。

$Q_k$ 使用真实目标梯度定义，含噪参数更新产生的实际位移则取决于梯度估计。比较同一目标下的估计器时，固定 $a_0$、$\gamma$、$\Theta$ 和目标设定；比较动态、静态参考点或不同效用方法时，目标本身随方法改变，残差数值须结合各自目标尺度解释。

\section{分析假设}\label{subsec:assumptions}

梯度估计与在线跟踪依赖不同层次的正则性。首先，有界分析域和似然条件保证 CPT 梯度存在，固定正则化控制有限样本估计中的非线性余项。其次，目标变化与模拟误差决定在线上界中的时变项和统计项。参考点敏感性、驻点距离及极小极大精度则分别引入额外的耦合、几何和观测条件。

以下各常数对所分析的时期、参数以及真实模型和模拟模型一致成立。关于 $K\to\infty$ 的陈述还要求这些常数对所有 $K$ 保持统一。模型定义和解析界用于确立这些条件，有限样本诊断用于考察实际运行是否接近相应边界。

\begin{assumption}[有界定义域与似然正则性]\label{ass:bounded}
参数集 $\Theta\subset\R^d$ 非空、紧且凸，并包含于一个定义策略的开参数邻域。特征可测，不显式依赖 $\theta$，且满足 $\norm{q_{i,u}}_2\le Q_f$。时域长度和资产数量均有限且固定。初始状态以及所分析的全部真实轨迹和模拟轨迹满足统一界 $0<X_u\le B_X$、$0\le r_u\le B_r$ 与 $0\le v_\pm(Y_k)\le B_v$。似然计算使用 $\Theta$ 的一个有界开邻域 $O$，其闭包包含于参数定义域。求导时固定初始信息，并且轨迹分布对 $\theta$ 的唯一显式依赖来自式~\eqref{eq:dirichlet-policy} 的 Dirichlet 动作密度；条件市场转移核与参数无关。引理~\ref{lem:dirichlet-moments} 从这些原始条件出发，证明轨迹积分下两次求导的可积支配条件以及一致的固定阶得分矩。将其中的得分二阶矩界与对数似然 Hessian 的期望界分别记为 $C_{G,2}$ 和 $C_H$。二者是推导得到的有限常数，并非得分函数的逐点上界。
\end{assumption}

该条件为似然比求导与可积性提供依据。参数集紧性和特征有界性可以从定义直接检查，财富界与价值界则来自给定市场模型、有界定义域或经过分析的停止规则。样本得分矩能够显示极端权重对数值计算的影响，总体矩界仍由后文的似然分析建立。若实际模型超出这一分析域，使用梯度恒等式与有限方差结论需要重新验证支配条件；有限回放所得到的界还须考察其对总时长的依赖。

\begin{assumption}[固定的正则化 CPT]\label{ass:smooth-cpt}
归一化权重函数递增且属于 $C^3([0,1])$，并满足
\[
 C_j:=\max_{\pm}\sup_{p\in[0,1]}|(w_\pm^\eps)^{(j)}(p)|<\infty,
 \qquad j=1,2,3.
\]
所有价值参数和权重参数保持固定。价值函数严格采用式~\eqref{eq:value-functions} 的形式，并使用其中固定的正数 $\eps_v$。定理~\ref{thm:gradient-identity} 根据引理~\ref{lem:dirichlet-moments} 推导出共同的目标梯度 Lipschitz 界
$L=2B_v[C_2C_{G,2}+C_1(C_{G,2}+C_H)]$。因此，目标函数的光滑性来自原始似然条件。价值函数平滑用于获得关于参考点的线性敏感性，关于策略参数的求导则通过轨迹似然进行。
\end{assumption}

上述导数界用于证明 CPT 目标可微，并控制乘子 $f(p,I)=(w^\eps)'(p)(I-p)$ 的二阶余项。固定正则化区间上的端点值和各阶导数可以解析检查。实现时，不同层级必须使用相同的权重函数及其导数；数值 Hessian 检查只能作为诊断，不能代替统一光滑性证明。若正则化参数随样本量变化，目标本身也随之变化，后文针对同一固定目标的多层去偏推导便需要重新论证。端点处奇异且未经正则化的权重函数不属于该假设的范围。

\begin{assumption}[分别刻画目标变化与模拟误差]\label{ass:drift}
对 $k<K$，定义下列有限可测量：
\begin{align}
 \delta_k^{J}&=\sup_{\theta\in\Theta}|J_{k+1}(\theta,r_{k+1})-J_k(\theta,r_{k})|,
 &V_K&=\sum_{k=1}^{K-1}\delta_k^{J},\label{eq:value-drift}\\
 \delta_k^{g}&=\sup_{\theta\in\Theta}
 \norm{\nabla J_{k+1}(\theta,r_{k+1})-\nabla J_k(\theta,r_{k})}_2.&&\label{eq:gradient-drift}
\end{align}
在被查询参数处，模拟误差满足
\begin{equation}\label{eq:simulation-bias}
 \norm{\widetilde g_k-g_k}_2\le\beta_k^{\rm sim},
\end{equation}
其中 $\beta_k^{\rm sim}$ 为 $\cF_k$ 可测。每个界中使用的期望均有限。除非在某个渐近陈述中另行明确规定，否则不预设这些量足够小或关于 $K$ 次线性增长。
\end{assumption}

函数值变化用于望远镜求和，梯度变化用于驻点集合漂移，模拟偏差则进入梯度估计均方误差。在给定市场模型下，可以解析控制这些量；使用观测数据时，可以借助拟合模型差异和独立梯度估计对其进行估计，但这些估计也带有自身误差。若变化预算随 $K$ 线性增长，或模拟误差持续存在，理论上界中可能保留不消失的项。时期内的真实变化由 $P_k$ 描述，模拟器对这部分变化的遗漏通过 $\beta_k^{\rm sim}$ 评价。$\beta_k^{\rm sim}$ 衡量同一期两个模型之间的差异，$V_K$ 衡量真实条件目标的跨期差异。

\begin{assumption}[参考点隔离耦合]\label{ass:reference-coupling}
讨论参考点敏感性时，固定 $P,z,\theta$，只改变初始参考点，并要求策略特征与市场转移核不直接依赖该初始参考点或其后续递推值。因此，使用相同的策略随机数和市场随机数时，两条耦合轨迹具有相同的财富路径和投资组合路径。允许的初始参考点区间覆盖变化分解中使用的全部比较。
\end{assumption}

该条件允许在共享得分函数与财富路径的条件下，逐路径比较两条参考点递推。可以通过检查状态到特征的映射以及市场转移核来判断条件是否满足。例如，算法~\ref{alg:dirichlet} 所列因子可以不包含 $r$，同时学习参数仍可通过 CPT 梯度对参考点作出响应。若特征使用 $X-r$ 或直接使用 $r$，基本的估计器结论和残差定理仍然可用，但在没有新的联合耦合论证时，不再主张隔离参考点的敏感性结论及其显式变化分解成立。

\begin{assumption}[残差误差界，可选条件]\label{ass:error-bound}
对式~\eqref{eq:normalized-map} 中固定的 $a_0,\gamma,\Theta$，存在统一常数 $\kappa<\infty$，使得
\begin{equation}\label{eq:error-bound}
 \dist(\theta,\Sta_k)\le\kappa\norm{Q_k(\theta)}_2,
 \qquad \theta\in\Theta.
\end{equation}
\end{assumption}

这一附加几何条件把小残差转换为参数到驻点集合的距离界。它必须来自给定目标函数的数学性质。采用多个初始点的局部搜索可以探查该条件，但不能构成证明；如果只在某个邻域内成立，还需要证明迭代点与比较点始终位于该邻域。不具备该条件时，主要残差界与估计器结果仍然有效，只是不能进一步推出驻点集合距离界。例如，$-x^4$ 一类平坦目标在驻点附近可能不满足线性残差误差界。

\begin{assumption}[统计模型类与精确轨迹信息]\label{ass:oracle}
固定 $d\geq1$、内部目标参数 $\theta_0=0\in\Theta$、有限上界 $\bar h,\bar N\geq1$，以及假设~\ref{ass:bounded} 和假设~\ref{ass:smooth-cpt} 中的共同常数。固定本文的平滑价值函数和归一化权重函数，并要求 $(w_+^\eps)'(0)>0$。考虑条件投资组合模型类 $\mathfrak C$，其中模型满足 $1\le h\le\bar h$、$1\le N\le\bar N$ 以及上述共同界。时域长度、资产数量、特征、参考点参数和起始信息均已知，市场分布未知。选取的常数应覆盖附录~\ref{app:minimax-proof} 在固定市场概率区间内构造的单期模型；在该子模型中，参考点适应率、收益幅度和起始财富也均为已知量。

关于未知市场分布的全部信息均来自精确预言机。算法可以在任意选定的 $\theta\in\Theta$ 处调用预言机，每次调用返回该策略下的一条全新完整轨迹。在给定查询参数与模型的条件下，不同调用相互独立。估计器可以自适应选择查询参数，可以使用与模型无关的随机化，并可以在其观测历史对应的停止时刻 $N_{\rm call}$ 停止，但必须满足
\[
 \sup_{P\in\mathfrak C}\E_P N_{\rm call}\le B.
\]
模型不提供关于未知市场分布的免费观测、已知市场分布参数，或包含额外未知分布观测的拟合模型辅助信息。估计目标是共同欧氏参数坐标下的 $g_P=\nabla J_P(\theta_0)$。当 $B\to\infty$ 时，维度与全部定义域常数保持固定。
\end{assumption}

该假设确定了极小极大分析的观测实验。每条完整轨迹计为一次调用，未知市场分布的信息通过这些调用逐步获得；模型类和维度固定，使不同估计程序能够在同一期望预算下比较。权重函数的正端点斜率可以直接检查，并使后文两点子模型中的 CPT 梯度具有局部可识别性。端点斜率为零时，估计上界仍成立，下界则需要另行构造可识别的子模型。使用拟合模拟器、额外市场观测、确定性预算或增长维度，会改变这里的观测条件和比较范围，需要相应的统计分析。

前两项假设支撑梯度与估计器结果。目标变化和模拟精度进入在线优化界。参考点隔离条件与残差误差界分别用于敏感性结论和驻点集合距离结论。精确预言机假设只用于极小极大定理。这样的分工使每项结论依赖的前提保持清楚，也避免把理想化统计实验与实际拟合模拟器的误差混在同一条件中。
